\section*{Preface}
\addcontentsline{toc}{section}{Preface}

Robert A. Di Paola and Alex Heller sought an algebraic treatment of
generalized incompleteness. At the end of Section~8 of their 1987 paper
they named their diagonal inseparability theorem among the results
they expected to use \cite[Theorem~8.15 and p.~629]{DPH}. Their
announced treatment did not appear. This monograph supplies an
algebraization in recursion categories. Under additional hypotheses,
we construct internal first-order theories in the restriction idempotent
completion of a recursion category. A consistent theory freely generated
by a disjoint pair of restriction idempotents is essentially incomplete
exactly when the pair is inseparable. A comparison theorem for range
functors gives conditions under which the diagonal pair of a Turing
morphism is inseparable.

To make the construction internal, we use syntax with named variables,
define substitution when the term is free for the variable, and obtain
derivations from initial algebras. A disjoint pair on an object $X$ of
propositional atoms specifies which atoms to assert and which to refute.
Suppose the theory generated by the pair is consistent. In a consistent
target theory, assign a sentence to every atom of $X$ so that the target
proves the sentences assigned to the first family and refutes those
assigned to the second. The generated theory, with its atom constructor,
is initial among consistent theories equipped with such an assignment.
Each assignment extends uniquely to a translation from the generated theory.

A complete consistent extension would separate the pair. Shoenfield used
the corresponding inverse-image argument in his proof of essential
undecidability \cite[Theorem~2, pp.~390--391]{Shoenfield}. Conversely,
a separator gives a Boolean valuation whose true sentences form a
complete consistent extension. Each theory determines a recursion
category over its theorem object. An extension of theories induces a
fully faithful functor between the corresponding completions, both left
and right adjoint to pullback.

In particular, the pullback argument of
Corollary~\ref{cor:domain-incompleteness} gives a categorical form of
Smullyan's conclusion \cite[Chapter~III, Theorem~23(b)]{Smullyan}.
If a theory admits such an assignment on an inseparable pair, the
theorem and refutation subobjects of every consistent extension are
inseparable. In arithmetic, Corollary~\ref{cor:arithmetic} applies
the comparison theorem to the standard interpretation of $S_{\PA}$
in sets and partial functions. It follows that in every consistent
computably enumerable extension of Peano arithmetic the sets of
numerical codes of provable and refutable sentences are recursively
inseparable. The other two applications use
computation uniform in the binary representations of a Miller sequence
and cosheaves on a three-point space.

Palomar entry \texttt{PALOMAR-2026-11-22-000009}, version~3,
registers three selected declarations. Appendix~\ref{app:reading-lean-proofs}
gives the correspondence with the Lean formalization.

Codex wrote Lean proofs under my direction. ChatGPT assisted with the
mathematical development and manuscript revisions. Claude (Anthropic)
reviewed manuscript drafts and checked citations and attributions against
the sources. ChatGPT and Claude each corrected errors in the other's
suggestions. I am responsible for the manuscript and its formalization.
